% !TeX spellcheck = ru_RU
% !TEX root = vkr.tex

\section*{Заключение}
\label{sec:conclusion}

В ходе работы была разработана система для сегментации, контурирования и параметрического
анализа биологических объектов по изображениям.
Система реализована в виде кроссплатформенного настольного приложения на Python (PyQt6)
и не требует GPU для работы.

Получены следующие результаты:

\begin{enumerate}
    \item Проведён анализ предметной области и существующих инструментов.
          Показано, что ни одно из рассмотренных решений (ImageJ, tpsDIG2, PlantCV)
          не покрывает полный рабочий цикл биолога в едином приложении.

    \item Сформулированы функциональные и нефункциональные требования совместно с консультантом
          из лаборатории Гербарий ГБС РАН.

    \item Разработана архитектура системы, включающая модуль сегментации на основе MobileSAM,
          интерактивный холст с режимами подсказок и ручной коррекции,
          буфер масок и модуль пакетных измерений.

    \item Реализован алгоритм параметрических измерений на основе PCA,
          вычисляющий длину и поперечные профили ширины листа в семи точках вдоль главной оси.

    \item Проведено сравнение автоматической сегментации с ручной разметкой биологов
          на выборке из 9~листьев.
          % TODO: вставить итоговые числа IoU и отклонения после получения данных от биологов

    \item Система апробирована консультантом; качество сегментации признано удовлетворительным.
\end{enumerate}

Направления дальнейшего развития:
\begin{itemize}
    \item реализация масштабного коэффициента для вывода измерений в миллиметрах;
    \item расстановка точек вдоль периметра с заданным шагом;
    \item дообучение MobileSAM на специфических снимках Sphagnum для повышения точности;
    \item расширение тестовой выборки и проведение количественной апробации с биологами.
\end{itemize}
